\section{Characterizations of invertible matrices}
The Invertible Matrix Theorem (IMT) is a collection of equivalent tests for
an $n\times n$ matrix. Invertibility, $n$ pivots, independent columns,
columns spanning $\R^n$, a trivial null space, and a one-to-one and onto
map $x\mapsto Ax$ all describe the same situation. The word
\emph{square} is essential: a rectangular matrix can have independent
columns without spanning its whole codomain.

For implications, distinguish a proof from a plausible example. A true
implication needs an argument for \emph{every} allowed matrix; a false
implication needs just one counterexample with true premise and false
conclusion. The zero matrix and identity matrix are useful tests, but their
special properties should not be mistaken for universal properties.

\practice{2.3.P1}{Repeated columns prevent invertibility}
The three columns of the supplied matrix are $(2,2,2)^T$,
$(3,3,3)^T$, and $(4,4,4)^T$. In particular column 2 is $3/2$ times
column 1, so the columns are dependent. The IMT says an invertible square
matrix has independent columns. Therefore this matrix is not invertible.

\practice{2.3.P2}{Negate an existence statement correctly}
Statement (g) says: for \emph{every} $b\in\R^n$ there exists a solution
to $Ax=b$. Its negation is: there exists \emph{at least one} $b$ for which
$Ax=b$ is inconsistent. It does not say every system is inconsistent;
indeed $Ax=0$ always has $x=0$ as a solution. The IMT also implies that
this $A$ is not invertible.

\practice{2.3.P3}{Recognize a noninvertible product}
The matrix $AB$ is square. A nontrivial solution of $(AB)x=0$ violates
the trivial-null-space condition in the IMT, so $AB$ is not invertible.
We cannot conclude that \emph{both} $A$ and $B$ are singular; one singular
factor suffices to make a product singular. For instance $A=0$, $B=I$
has exactly one singular factor.

\exercise{2.3.1}{A two-by-two determinant test}
For $A=\begin{bmatrix}5&7\\-3&-6\end{bmatrix}$,
$ad-bc=-30+21=-9\ne0$. The two-by-two inverse criterion from Section 2.2
therefore gives \textbf{invertible}. No complete inverse calculation is needed.

\exercise{2.3.2}{A multiple row}
The second row $(6,-3)$ is $-3/2$ times the first row $(-4,2)$.
Eliminating it leaves only one pivot, fewer than two. Thus the matrix is
\textbf{not invertible}. Equivalently the determinant is $12-12=0$.

\exercise{2.3.3}{A triangular matrix with nonzero diagonal}
The matrix is lower triangular with diagonal $3,-4,-3$.
The first pivot $3$ removes the entries below it; the second pivot $-4$
removes the next entry below it; the last pivot remains $-3$. None vanishes.
There is a pivot in every row and column, so the IMT gives
\textbf{invertible}. The triangular structure saves a full reduction.

\exercise{2.3.4}{A zero row}
The second row is all zeros. Thus there can be at most two pivot rows in
this $3\times3$ matrix. Since invertibility would require three pivots,
the matrix is \textbf{not invertible}.

\exercise{2.3.5}{A zero column}
Column 2 is zero, so $Ae_2=0$ with $e_2\ne0$. This is a nontrivial
null vector. The IMT therefore gives \textbf{not invertible}; equivalently
the columns are dependent because one of them is the zero vector.

\exercise{2.3.6}{A short elimination gives the last pivot}
Use $R_3\gets R_3+3R_1$ to obtain the last row $(0,-3,-18)$,
then $R_3\gets R_3+(3/4)R_2$. The resulting echelon matrix is
\[
 \begin{bmatrix}1&-3&-6\\0&4&3\\0&0&-63/4\end{bmatrix}.
\]
Its three diagonal entries are nonzero pivots, so the original matrix is
\textbf{invertible}. The operations preserve row equivalence.

\exercise{2.3.7}{Four pivots from a few row replacements}
Apply $R_2\gets R_2+3R_1$, $R_3\gets R_3-2R_1$, and then
$R_4\gets R_4-(1/4)R_2$. This gives
\[
 \begin{bmatrix}-1&-3&0&1\\0&-4&8&0\\0&0&3&0\\0&0&0&1\end{bmatrix}.
\]
Every diagonal position is a pivot. Hence the original $4\times4$ matrix
is \textbf{invertible}. There is no need to clear the entries above the pivots.

\exercise{2.3.8}{Read pivots from an upper triangular matrix}
The given matrix is already in echelon form with nonzero diagonal entries
$3,1,2,1$. It has four pivots, so it is \textbf{invertible} by the IMT.
The entries above the diagonal do not change this conclusion.

\exercise{2.3.9}{A computer-assisted exact pivot test}
Using rational row reduction on the supplied integer matrix gives an
echelon matrix
\[
 \begin{bmatrix}
 4&0&-3&-7\\0&9&9/2&-3/2\\0&0&71/4&365/12\\0&0&0&-1336/213
 \end{bmatrix}.
\]
For reproducibility, elimination here uses
$R_i\gets R_i-(a_{ik}/a_{kk})R_k$ below each successive pivot.
All four pivots are nonzero; therefore the matrix is \textbf{invertible}.
An exact program check is \texttt{Matrix(A).rref()}, using SymPy, which
returns $I_4$ and pivot columns $(0,1,2,3)$ (zero-based indices).
These exact calculations avoid confusing a very small pivot with zero.

\exercise{2.3.10}{Five pivots despite a small determinant}
For the supplied $5\times5$ integer matrix, exact elimination, including
one row interchange when a candidate pivot is zero, gives
\[
 \begin{bmatrix}
 5&3&1&7&9\\0&2/5&4/5&-2/5&-94/5\\
 0&0&1&-21&7\\0&0&0&1&34\\0&0&0&0&-1
 \end{bmatrix}.
\]
There are five nonzero pivots, hence the matrix is \textbf{invertible}.
Its inverse can nevertheless be numerically sensitive; invertibility is
an exact yes/no statement, while conditioning measures sensitivity.
Exercise 43 quantifies that distinction. Exact SymPy row reduction returns
$I_5$ with all five pivot columns.

\exercise{2.3.11}{Five implications, with the square hypothesis retained}
\begin{enumerate}[label=(\alph*)]
\item \textbf{True.} A trivial null space means no free variables,
so all $n$ columns have pivots. For a square matrix that forces the reduced
form to be $I_n$; this is an IMT equivalence.
\item \textbf{True.} Spanning $\R^n$ forces $n$ pivot rows. In an
$n\times n$ matrix these are also $n$ pivot columns, giving independence.
\item \textbf{False.} Merely being $n\times n$ does not guarantee
consistency for every right side. Take $A=0$ and $b=e_1$; the equation
$0x=e_1$ has no solution.
\item \textbf{True.} A nontrivial homogeneous solution requires a free
variable, so some column is not a pivot column. Thus there are fewer than
$n$ pivots.
\item \textbf{True.} If $A$ were invertible, then
$(A^{-1})^T$ would be an inverse of $A^T$. The given noninvertibility of
$A^T$ excludes that possibility, so $A$ is not invertible.
\end{enumerate}

\exercise{2.3.12}{Pay attention to onto versus into}
\begin{enumerate}[label=(\alph*)]
\item \textbf{True.} $AD=I$ is the right-inverse condition in the IMT,
so $A$ is invertible. Left multiplication by $A^{-1}$ gives $D=A^{-1}$,
and therefore $DA=I$.
\item \textbf{False.} Every $n\times n$ matrix maps $\R^n$ \emph{into}
$\R^n$, including $A=0$, whose reduced form is not $I$.
Mapping \emph{onto} would be a stronger premise and would imply the conclusion.
\item \textbf{True.} For a square matrix, independent columns and
columns spanning $\R^n$ are equivalent by the IMT.
\item \textbf{False.} Consistency for every $b$ implies invertibility,
hence the transformation \emph{is} one-to-one. The identity matrix is a
concrete counterexample to the claimed negative conclusion.
\item \textbf{False.} Consistency for just one right side says nothing
about uniqueness. For $A=0$, $b=0$, every vector is a solution.
\end{enumerate}

\exercise{2.3.13}{When an upper triangular matrix is invertible}
It is invertible \emph{if and only if every diagonal entry is nonzero}.
If all are nonzero, eliminate above the diagonal from the last row upward,
then scale the rows; this reduces the matrix to $I_n$.
For the converse, suppose diagonal entry $k$ is zero. In an upper
triangular matrix, each of the first $k-1$ columns has zero entries from
row $k$ downward. Column $k$ also has zero entries from row $k$ downward:
its diagonal is zero, and its lower entries vanish by triangularity.
Thus the first $k$ columns live entirely in their first $k-1$ coordinates.
A homogeneous system with $k$ variables but only $k-1$ potentially
nonzero equations has a free variable, so these $k$ columns are dependent.
Their dependence extends to all columns by assigning zero coefficients
to the remaining columns. The IMT therefore excludes invertibility.

\exercise{2.3.14}{The lower triangular case}
Again the criterion is that every diagonal entry be nonzero.
One way to avoid repeating the argument is to transpose: $A$ is lower
triangular exactly when $A^T$ is upper triangular, with the same diagonal.
The IMT says $A$ is invertible exactly when $A^T$ is invertible.
Applying Exercise 13 to $A^T$ proves both directions of the claimed criterion.

\exercise{2.3.15}{Nonspanning columns exclude an inverse}
No. The IMT says an invertible $4\times4$ matrix has columns spanning
$\R^4$. The proposed failure of spanning is the negation of one of its
equivalent conditions, so it rules out invertibility.

\exercise{2.3.16}{Independence of the transposed columns}
If $A$ is invertible, then $A^T$ is invertible, with inverse $(A^{-1})^T$.
Applying the independent-column condition of the IMT to $A^T$ shows that
its columns are independent. Equivalently the rows of $A$ are independent.

\exercise{2.3.17}{Identical columns produce an explicit dependence}
No. If columns $i,j$ are equal with $i\ne j$, then
$A(e_i-e_j)=a_i-a_j=0$. The vector $e_i-e_j$ is nonzero, so the null
space is nontrivial and $A$ cannot be invertible.

\exercise{2.3.18}{Identical rows also rule out invertibility}
No. Equal rows of $A$ become equal columns of $A^T$.
Exercise 17 makes $A^T$ singular, and the transpose equivalence in the
IMT then makes $A$ singular too.

\exercise{2.3.19}{Seven independent columns give every right side exactly once}
Since $D$ is $7\times7$ and has independent columns, the IMT makes it
invertible. For \emph{every} $b\in\R^7$, the system $Dx=b$ therefore
has the unique solution $x=D^{-1}b$. Both existence and uniqueness are
part of the conclusion; neither needs a special choice of $b$.

\exercise{2.3.20}{Universal consistency for a square system implies uniqueness}
It is not possible to have more than one solution for any $b$.
Universal consistency gives a pivot in every one of the five rows.
Because there are five columns too, all five variables are pivot variables,
leaving no free variable. The matrix is invertible, and every system has
exactly one solution.

\exercise{2.3.21}{Multiple solutions exclude spanning}
If $Cu=v$ has two different solutions $u_1,u_2$, then
$C(u_1-u_2)=v-v=0$ with $u_1-u_2\ne0$. Thus $C$ is singular by the
IMT. Its columns therefore cannot span $\R^n$. Notice that the original
right side $v$ is still attained; the failure concerns some other right sides.

\exercise{2.3.22}{A one-sided square inverse forces commutation}
From $EF=I$, the IMT makes $E$ invertible. Multiplying on the left by
$E^{-1}$ gives $F=E^{-1}$, so $FE=I$ as well.
Hence $EF=FE$, which is exactly the definition of commuting matrices.
This conclusion uses squareness; rectangular one-sided inverses need not reverse.

\exercise{2.3.23}{A missed output implies a nontrivial null vector}
If $Fx=y$ is inconsistent for some $y$, then $F$ is not onto.
By the IMT this square matrix is singular, so $Fx=0$ has a nontrivial
solution. In fact if $z\ne0$ is such a solution, every $tz$, $t\in\R$,
is a solution; thus the homogeneous system has infinitely many solutions.

\exercise{2.3.24}{A failure to reduce to the identity}
If $G$ cannot be reduced to $I_n$, the IMT makes it singular.
Consequently its columns are dependent and do not span $\R^n$.
There are fewer than $n$ pivots. These are equivalent descriptions of the
same failure, not alternative cases to choose between.

\exercise{2.3.25}{Verify the boxed fact about a one-sided inverse}
Assume square matrices $A,B$ satisfy $AB=I$. Then $A$ has a right
inverse, so the IMT implies that $A$ is invertible. The equation gives
$B=A^{-1}AB=A^{-1}$. Therefore $B$ is invertible too, with inverse $A$,
and $BA=I$. This proves every assertion in the boxed statement preceding
Example 1, rather than merely assuming a right inverse is automatically two-sided.

\exercise{2.3.26}{Independent columns lead to a spanning square}
Independence of the columns of square $A$ makes $A$ invertible.
The product $A^2=AA$ is then invertible, with inverse $A^{-1}A^{-1}$.
Applying the spanning condition of the IMT to $A^2$ proves that its columns
span $\R^n$. Equivalently, $A$ maps onto $\R^n$ and doing so twice
still reaches every vector.

\exercise{2.3.27}{An invertible product forces the first factor invertible}
Let $W=(AB)^{-1}$. Then $A(BW)=(AB)W=I$.
Thus $A$ has a square right inverse $BW$, and the IMT makes $A$
invertible. This proof does \emph{not} apply the inverse-of-a-product
formula to unverified inverses of $A,B$; their existence is what must be proved.

\exercise{2.3.28}{An invertible product forces the second factor invertible}
Suppose $Bx=0$. Multiplying by $A$ gives $(AB)x=0$.
Invertibility of $AB$ implies $x=0$. Hence $B$ has trivial null space,
and since $B$ is square the IMT proves it is invertible.
Alternatively, $(WA)B=I$ with $W=(AB)^{-1}$ gives a left-inverse test.

\exercise{2.3.29}{One-to-one endomorphisms are also onto}
The one-to-one condition says $Ax=0$ has only $x=0$, or equivalently the
columns are independent. For $n\times n$ matrices the IMT makes $A$
invertible and hence onto $\R^n$. Thus the transformation is a bijection
with inverse $x\mapsto A^{-1}x$. The equal dimensions are crucial.

\exercise{2.3.30}{Failure of injectivity also forces failure of surjectivity}
Two solutions of $Ax=b$ give a nonzero homogeneous solution by subtraction.
Therefore $A$ is singular. By the IMT the associated transformation is
not onto $\R^n$ either: at least one right side has no preimage.
It can still attain the particular $b$ mentioned in the premise.

\exercise{2.3.31}{Prove uniqueness directly from pivot counting}
\goal Use row reduction, without Theorems 5 or 8.
Consistency for all $b$ forces a pivot in every row; otherwise an
appropriate right side creates an inconsistent zero-coefficient row.
The square matrix has $n$ rows and $n$ columns, so these $n$ pivots use
all $n$ columns. There are no free variables. Thus, for any fixed $b$,
back-substitution determines each variable uniquely. The hypothesis gives
existence, and this pivot argument gives uniqueness, as requested.

\exercise{2.3.32}{Prove existence directly from a trivial null space}
\goal Avoid invoking the IMT.
If a column were not a pivot column, its variable could be chosen freely
in the homogeneous system, producing a nonzero solution. Thus all $n$
columns have pivots. The $n$ pivots occupy different rows, and there are
only $n$ rows, so every row has a pivot. Hence for every augmented right
side the system is consistent: no zero coefficient row can create a
contradiction. This gives a solution of $Ax=b$ for every $b$ directly.

\exercise{2.3.33}{Invert a two-dimensional linear transformation}
The standard matrix is
$A=\begin{bmatrix}-5&9\\4&-7\end{bmatrix}$, whose determinant is
$35-36=-1$. Thus
$A^{-1}=\begin{bmatrix}7&9\\4&5\end{bmatrix}$.
For an output $(y_1,y_2)$, its unique input is
\[
 \boxed{T^{-1}(y_1,y_2)=(7y_1+9y_2,\ 4y_1+5y_2)}.
\]
Multiplying either pair of matrices gives $I_2$, which verifies both
compositions $T^{-1}\circ T$ and $T\circ T^{-1}$ are the identity.

\exercise{2.3.34}{A second inverse transformation}
Now $A=\begin{bmatrix}2&-8\\-2&7\end{bmatrix}$ has determinant
$14-16=-2$, so
$A^{-1}=\begin{bmatrix}-7/2&-4\\-1&-1\end{bmatrix}$.
Consequently
\[
 \boxed{T^{-1}(y_1,y_2)=(-\tfrac72y_1-4y_2,\ -y_1-y_2)}.
\]
For example applying $T$ to this pair gives first coordinate
$-7y_1-8y_2+8y_1+8y_2=y_1$ and second coordinate
$7y_1+8y_2-7y_1-7y_2=y_2$. The reverse composition follows from the
same inverse matrix calculation.

\exercise{2.3.35}{Use both inverse identities directly}
Let $S=T^{-1}$, so $S(T(x))=x$ and $T(S(y))=y$.
If $T(u)=T(v)$, apply $S$ to get $u=v$; thus $T$ is one-to-one.
For every desired output $y$, the input $S(y)$ satisfies $T(S(y))=y$;
thus $T$ is onto. As a second proof, Theorem 9 identifies invertibility
of $T$ with invertibility of its square standard matrix, and the IMT
then gives both injectivity and surjectivity.

\exercise{2.3.36}{Two inputs sharing an output rule out onto}
The equality $T(u)=T(v)$ for $u\ne v$ means $T$ is not one-to-one.
Its square standard matrix therefore fails one IMT condition, so it also
fails the onto condition. Hence $T$ cannot map $\R^n$ onto $\R^n$.
One can also subtract to see $T(u-v)=0$ with $u-v\ne0$.

\exercise{2.3.37}{A one-sided linear inverse in equal dimensions}
Let $A,B$ be the standard matrices of $T,U$. The hypothesis
$T(U(x))=x$ for all $x$ implies $ABx=x$ for all $x$.
Testing the standard basis vectors gives $AB=I$.
The boxed result proved in Exercise 25 yields $BA=I$.
Thus $U(T(x))=BAx=x$ for every $x$. The conclusion depends on both
linearity and the common finite dimension.

\exercise{2.3.38}{Why the inverse function is unique}
Since $T$ is invertible, it is onto. Given arbitrary $v\in\R^n$,
choose $x$ with $T(x)=v$. Then
\[
 S(v)=S(T(x))=x=U(T(x))=U(v).
\]
Every $v$ was allowed, so $S=U$ as functions. We did not need to assume
that $S,U$ were linear; the surjectivity of $T$ ensured their arguments
covered the entire domain.

\exercise{2.3.39}{An onto linear map in equal dimensions has an onto inverse}
By the IMT applied to its standard $n\times n$ matrix, an onto $T$
is invertible. For any $x$, taking $y=T(x)$ gives $T^{-1}(y)=x$;
therefore $T^{-1}$ is onto. If $T^{-1}(y_1)=T^{-1}(y_2)$,
apply $T$ to obtain $y_1=y_2$; therefore it is one-to-one as well.
So both $T$ and its inverse are bijections of $\R^n$.

\exercise{2.3.40}{Prove linearity of the inverse without assuming it}
Let $x=S(u)$ and $y=S(v)$. The identity $T\circ S=\mathrm{id}$ gives
$T(x)=u$, $T(y)=v$. By linearity of $T$,
$T(x+y)=u+v$. Applying $S$ and using $S\circ T=\mathrm{id}$ yields
\[
 S(u+v)=S(T(x+y))=x+y=S(u)+S(v).
\]
Likewise $T(cx)=cT(x)=cu$ gives
$S(cu)=S(T(cx))=cx=cS(u)$. These are additivity and homogeneity,
so $S$ is linear. Both inverse identities were used explicitly.

\exercise{2.3.41}{Tiny data perturbations, large solution changes}
\step{(a) Solve both systems exactly.}
For either right side $(b_1,b_2)$, the inverse of the coefficient matrix is
\[
 A^{-1}=\frac1{4.5(1.1)-3.1(1.6)}
 \begin{bmatrix}1.1&-3.1\\-1.6&4.5\end{bmatrix}
 =\begin{bmatrix}-110&310\\160&-450\end{bmatrix},
\]
because $4.95-4.96=-.01$. For the unrounded right side,
\[
 x_1=-110(19.249)+310(6.843)=3.94,\qquad
 x_2=160(19.249)-450(6.843)=.49.
\]
For the rounded right side, the same products give $(2.90,2.00)^T$.
All these decimals are exact terminating decimals, not numerical approximations.
\step{(b) Specify the meaning of percentage error.}
Coordinatewise relative percentage errors are
\[
 100\frac{|2.90-3.94|}{3.94}=\frac{5200}{197}\%\approx26.40\%,
 \qquad
 100\frac{|2.00-.49|}{.49}=\frac{15100}{49}\%\approx308.16\%.
\]
If a single vector error is wanted instead, the Euclidean relative error is
$100\|\widetilde x-x\|_2/\|x\|_2\approx46.18\%$.
Stating the norm matters: these are different legitimate summaries.
\step{(c) Compute and interpret conditioning.}
Executing \texttt{np.linalg.cond(A)} in NumPy gives the spectral
($2$-norm) condition number $\kappa_2(A)\approx3362.9997$.
The $1$-norm and infinity-norm condition numbers both equal $4636$;
software conventions therefore need to be identified.
The large value explains why very small changes of the measured right
side can cause large changes in the recovered solution.

\exercise{2.3.42}{Measure roundoff for the matrix in Exercise 9}
\goal Compare a known vector with a solution recovered from its computed output.
The following is an executed NumPy experiment using binary64 arithmetic:
\begin{verbatim}
import numpy as np
A = np.array([[4,0,-3,-7],[-6,9,9,9],
              [7,-5,10,19],[-1,2,4,-1]],dtype=float)
x = np.random.default_rng(23042).random(4)
b = A@x
x1 = np.linalg.solve(A,b)
print(np.linalg.cond(A))
print(np.max(np.abs(x1-x))/np.max(np.abs(x)))
print(np.abs((x1-x)/x))
\end{verbatim}
The condition number is $11.3525089$, and the normwise relative error
$\|x_1-x\|_\infty/\|x\|_\infty$ is approximately $4.52\cdot10^{-16}$.
The chosen vector, rounded here for readability, is
\[
 x\approx(.98344656,.12402037,.75512517,.48561218)^T.
\]
The worst coordinatewise relative error is about $2.13\cdot10^{-15}$;
one coordinate was recovered exactly in this run.

Binary64 has $53$ binary significant bits, about $15$--$16$ decimal
significant digits. Thus the observed normwise agreement is roughly
$15$ digits, with roughly $14$--$15$ in the most sensitive coordinate.
The expected worst-case loss suggested by $\log_{10}\kappa_2(A)$ is about
$1.06$ digits. The measured error is an observation for this vector,
not a promise that every possible right side has identical accuracy.
Libraries and hardware may alter the last few floating-point bits.

\exercise{2.3.43}{Repeat with a more sensitive matrix}
Use the same procedure but with the matrix in Exercise 10 and the seed below:
\begin{verbatim}
A = np.array([[5,3,1,7,9],[6,4,2,8,-8],
              [7,5,3,10,9],[9,6,4,-9,-5],
              [8,5,2,11,4]],dtype=float)
x = np.random.default_rng(23043).random(5)
b = A@x
x1 = np.linalg.solve(A,b)
print(np.linalg.cond(A))
print(np.max(np.abs(x1-x))/np.max(np.abs(x)))
\end{verbatim}
The executed condition number is $\kappa_2(A)\approx68622.4548$;
the normwise relative error is approximately $6.01\cdot10^{-12}$.
Thus this run retains about $11$ decimal digits normwise, losing about
$4$--$5$ of the approximately $15$--$16$ available digits.
The theoretical sensitivity indicator $\log_{10}\kappa_2(A)\approx4.84$
is consistent with that behavior. The exact pivot calculation in Exercise 10
proved invertibility, but it did not guarantee numerical robustness.

\exercise{2.3.44}{The last inverse column of a fifth-order Hilbert matrix}
To obtain column 5 of $H^{-1}$ without the other columns, solve
$Hx=e_5$, where $H_{ij}=1/(i+j-1)$ for $1\le i,j\le5$.
An exact rational solve gives
\[
 \boxed{x=(630,-12600,56700,-88200,44100)^T}.
\]
It can be verified by multiplying the five rows of $H$ by this column;
the results are $0,0,0,0,1$ respectively. For example the first is
$630-6300+18900-22050+8820=0$.
A numerical version is
\begin{verbatim}
H = np.array([[1/(i+j+1) for j in range(5)]
              for i in range(5)])
x = np.linalg.solve(H,np.eye(5)[:,4])
print(np.linalg.cond(H))
\end{verbatim}
Here $\kappa_2(H)\approx476607.25\approx10^{5.68}$.
Starting with about $15$--$16$ decimal digits, the condition estimate
allows a loss of about six digits, so about $9$--$10$ digits is a
conservative expectation. Our executed solve has largest relative
coordinate error about $1.93\cdot10^{-12}$ against the exact column,
roughly $11$--$12$ digits of agreement. A bound on possible loss need
not be reached for this particular right side.

\exercise{2.3.45}{A twelfth-order Hilbert inverse can be numerically unreliable}
Use a numerical inverse, rather than an exact rational one, for this experiment:
\begin{verbatim}
H = np.array([[1/(i+j+1) for j in range(12)]
              for i in range(12)])
B = np.linalg.inv(H)
R = H@B-np.eye(12)
print(np.linalg.cond(H))
print(np.max(np.abs(R)))
print(np.linalg.norm(R,np.inf))
\end{verbatim}
In the executed NumPy binary64 run, the outputs are approximately
$1.64\cdot10^{16}$, $0.0892$, and $0.320$, respectively.
Thus the largest entry of the computed product's error is nearly $0.09$:
$HB$ is visibly different from the identity. The computed $B_{11}$ is
approximately $142.655$, whereas the exact Hilbert inverse has
$(H^{-1})_{11}=12^2=144$.

The exact rational Hilbert matrix is invertible, but its condition number
is comparable to the reciprocal of binary64 rounding precision. Ordinary
floating-point arithmetic therefore has essentially no reliable significant
digits guaranteed in an inverse computation of this size. Different
numerical libraries may produce somewhat different inaccurate answers;
none of that contradicts exact invertibility. Using higher precision and
checking against exact rational arithmetic are ways to investigate further,
not reasons to silently report the floating-point product as $I$.
